This paper speaks most directly to the economic geography of creative industries and to the small
literature on music scenes as localized production systems. \citet{scott_1999} provides the
classic economic-geography statement that recorded music production is concentrated in dense local
environments with specialized workers, organizations, and complementary institutions. The closest
empirical papers are \citet{klement_strambach_eg_2019} and
\citet{klement_strambach_rs_2019}. They show that related variety matters for urban music-scene
creativity and that new genres emerge from diversification processes that remain strongly local
even when extra-local influences matter. My paper fits squarely in that tradition, but shifts the
evidence base away from case studies and retrospective scene histories toward a large
city-by-genre panel built from linked band and member data.

The paper also contributes to a broader economics literature on agglomeration, local externalities,
and the formation of new varieties. \citet{duranton_puga_2004} provide the core micro-foundations:
dense places create advantages through sharing, matching, and learning. \citet{glaeser_etal_1992}
link local variety and dynamic externalities to urban growth, while \citet{storper_venables_2004}
stress the continuing value of face-to-face contact and local buzz even when communication
technologies improve. \citet{moretti_2004} provides direct evidence that local stocks of skilled
workers generate spillovers. These papers are not about music, but they provide the right
economics language for the paper's central claim: local scene thickness can lower the cost of
forming new creative projects because matching is easier, tacit know-how is more available, and
the local capability base is deeper.

The closest mechanism papers are on worker mobility, spinouts, and recombination.
\citet{franco_filson_2006} study spinouts and employee mobility as channels through which
organizational capabilities propagate, \citet{moen_2000} asks whether technical personnel mobility
is a source of R\&D spillovers, and \citet{weitzman_1998} provides the cleanest formal logic for
economically meaningful recombination. Bands are not firms in the standard sense, but the mapping
is useful: workers accumulate project-specific know-how, some leave to found new projects, and
overlapping experience across projects expands the scope for productive recombination. The paper's
spawning-flow and multi-band-musician variables are reduced-form traces of those mechanisms.

At a more abstract level, the paper uses a product-variety interpretation of genre emergence.
\citet{dixit_stiglitz_1977} and \citet{romer_1990} show why the creation of new varieties can be
economically meaningful even when the relevant outcome is not technological improvement in the
narrow sense. The point is not that metal scenes literally replicate aggregate growth theory. It is
that a new niche can be understood as a differentiated product environment that becomes locally
viable once enough project formation, capability, and specialization accumulate.

The paper is also informed by the economics of music markets under digitization.
\citet{ferreira_waldfogel_2013} show that geography and language continue to structure world music
trade, while \citet{waldfogel_2015} shows that digitization changed the supply side of recorded
music by lowering production and distribution frictions and expanding entry. These papers matter
mainly as boundary conditions. The question here is not whether digital technology eliminated
geography in music, but whether some places remain better at generating and sustaining local niche
production even in a more digital environment.

Finally, the paper's outcome is informed by work on genres as social forms rather than fixed
technical taxonomies. \citet{lena_peterson_2008} is especially important here. Genres are social
classifications with trajectories, organizational supports, and shifting boundaries, not objective
technologies with a single uncontested birth date. For that reason, the paper does not claim to
date the true global origin of a genre. Instead, it studies when a local city-genre niche becomes
thick enough to support repeated entry and collaboration. The empirical motivation in the next
section shows those forces in the raw data, and the model that follows compresses them into a
deliberately narrow object: a threshold-crossing entry problem driven by local niche thickness,
embedded spinouts, and recombinant worker depth.
